\documentclass[11pt]{article}
\usepackage[T1]{fontenc}
\usepackage[utf8]{inputenc}
\usepackage{lmodern,seqsplit,xurl}
\usepackage[margin=1in]{geometry}
\usepackage{amsmath,amssymb,amsthm}
\usepackage{array,booktabs,longtable,tabularx}
\usepackage{microtype}
\usepackage[hidelinks]{hyperref}
\newcolumntype{L}[1]{>{\raggedright\arraybackslash}p{#1}}
\newcommand{\code}[1]{\texttt{#1}}
\newcommand{\longcode}[1]{\texttt{\seqsplit{#1}}}
\emergencystretch=2em
\newcommand{\ty}[1]{\mathsf{#1}}
\newcommand{\ok}{\mathsf{ok}}
\newcommand{\some}{\mathsf{some}}
\newcommand{\rep}{\mathrel{\triangleright}}
\newcommand{\Nat}{\mathbb{N}}
\newcommand{\B}{\mathcal{B}}
\newcommand{\I}{\mathcal{I}}
\newcommand{\Pp}{\mathcal{P}}
\newcommand{\Ap}{\mathcal{A}}
\newtheorem{definition}{Definition}[section]
\newtheorem{proposition}[definition]{Proposition}
\hypersetup{
  pdftitle={The LeanExe Fragment: Types, Extraction, and Execution},
  pdfauthor={Codex GPT-6, Jamie Stephens},
  pdfsubject={Programming languages; implementation-indexed language specification}
}
\title{The LeanExe Fragment: Types, Extraction, and Execution}
\author{Codex GPT-6, Jamie Stephens}
\date{11 September 2026}

\begin{document}
\maketitle
\begin{abstract}
LeanExe compiles checked Lean declarations to WebAssembly through a restricted first-order extraction path.  The LeanExe fragment is specified here by runtime typing and layout rules together with an implementation-indexed relation from checked declarations to WebAssembly execution.  The report separates kernel typing, static specialization, extraction acceptance, ownership analysis, binary validity, and artifact behavior.  It defines the internal and public type domains, proves two layout propositions, and states the numeric, collection, memory, and fixed-adapter semantics used by the implementation.  A source audit identifies incomplete runtime export-name rejection and a child-mask width obligation.  The reference IR evaluator has a narrower comparison domain than compiled execution.  Existing emitter, decoder, validator, and artifact theorems support their named boundaries.  An independent source operational semantics, mechanized acceptance judgments, and a general source-to-WebAssembly refinement theorem remain open.
\end{abstract}

\section{Scope and evidence}

The LeanExe fragment is specified here by runtime typing and layout rules together with an implementation-indexed relation from checked declarations to WebAssembly execution.  An implementation-indexed relation is the graph of named compiler functions at a fixed source revision, composed with a stated target execution relation.  The independently stated typing rules describe runtime values and first-order code.  Shape-sensitive extraction premises still refer to compiler recognizers, so the report's acceptance definition depends on the implementation.

The source for this account is the specification checkpoint \longcode{2f3ec33f6e98ad98ac94df277c91e5319763cf9d} of LeanExe~\cite{checkpoint}.  The compiler root pins Lean 4.34.0-rc2 at commit \longcode{6a10ac8c22beadecabdbb0919c2b50214762f91d}.  The proof workspace pins Talos revision \longcode{87e3aa5e8f6e6f3b3eb5e7e4c5aba43071002d47}.  Source comparison at report preparation found no changes to the two specification documents or the inspected extraction, IR, and binary-emission files between that checkpoint and \longcode{7c7277d1e1ab8bc4ee5b4a014a8b975abb8b8054}.  The citations identify the fixed checkpoint throughout.

The evidence consists of source inspection, the mathematical arguments in this report, and existing theorem declarations cited with their exact scope.  Report preparation ran no Lean build, compiler, artifact verifier, or runtime reproducer.  In particular, the two implementation concerns in Section~\ref{sec:audit} have source evidence and remain untested as compiled examples.  The maintained language documents contain a fuller source-operation inventory~\cite{typetheory,specification}.  This report consolidates their mathematical objects and assurance boundaries.

Let $E$ be a checked Lean environment, $\Gamma$ a Lean local context, $\Delta$ a runtime-variable context, and $\Xi$ a first-order function-signature context.  The judgment $E;\Gamma\vdash_{\rm Lean} e:A$ is the pinned Lean kernel judgment.  Lean supplies dependent functions, universes, inductive formation, recursors, and definitional equality~\cite{lean}.  LeanExe reads elaborated declaration expressions from $E$, including static arguments and proof-bearing fields.

Kernel typing and runtime acceptance have separate premises.  A source theorem may quantify over arbitrary Lean types and unbounded natural numbers.  The executable declaration must also satisfy LeanExe's safety, body-availability, type, specialization, recursion, dependency, and ownership checks.  The default \code{ConstantInfo.value?} exposes definition bodies while excluding theorem and opaque bodies.  The executable signature checks reject unsafe and partial declarations~\cite{types,extract}.

\section{Runtime type discipline}

\subsection{Recognition and field erasure}

Runtime types have the following finite syntax:
\begin{align*}
\tau ::= {}&\ty{Unit}\mid\ty{Bool}\mid\ty{U8}\mid\ty{U32}\mid\ty{U64}
       \mid\ty{Nat64}\mid\ty{Bytes}\\
 &\mid\ty{Array}\;\tau\mid\tau\times\tau\mid\ty{PSum}\;\tau\;\tau\\
 &\mid\ty{Struct}(S,\bar\tau,\bar\upsilon)
  \mid\ty{Variant}(D,\bar\tau,[\bar\upsilon_0;\ldots;\bar\upsilon_{k-1}])
  \mid\ty{Rec}(D,\bar\tau).
\end{align*}
The names $S$ and $D$ identify Lean declarations.  The list $\bar\tau$ identifies concrete type parameters, and the lists $\bar\upsilon_i$ give retained fields.  The symbol $\ty{Nat64}$ denotes the bounded runtime interpretation of Lean \code{Nat}.  Recognized \code{PUnit} sequencing uses $\ty{Unit}$.  The abbreviations
\[
\ty{Option}\;\tau=\ty{Variant}(\code{Option},[\tau],[[];[\tau]])
\]
and
\[
\ty{Except}\;\epsilon\;\tau
=\ty{Variant}(\code{Except},[\epsilon,\tau],[[\epsilon];[\tau]])
\]
fix the constructor order and payload areas.

Fix an implementation revision $v$.  The implemented type-recognition and erasure judgments are
\begin{align*}
E\vdash_v A\rep\tau
&\quad\Longleftrightarrow\quad
\code{typeAtom?}(E,A)=\some(\tau),\\
E\vdash_v A\;\mathsf{erase}
&\quad\Longleftrightarrow\quad
\code{isProofType?}(E,A)=\mathsf{true}.
\end{align*}
The named functions occur in \code{Extract/Types.lean}~\cite{types}.  These judgments preserve the recognizers' syntactic limits.  A proof of kernel-definitional equality between two types alone does not establish equal recognizer results.

For a constructor field list, erase a field when the erasure judgment holds.  Otherwise retain its recognized runtime type, or reject when recognition fails.  The erasure test takes priority.  The implemented proof recognizer consumes metadata, recognizes \code{Sort 0}, follows the result of a dependent function type, and recognizes sufficiently applied declaration heads with proposition-valued declared results.  Completeness for all kernel propositions requires a separate proof.

A structure application must have concrete recognized parameters, no indices, one constructor, and a nonrecursive structure declaration.  The recognizer excludes classes and evidence-carrier names, substitutes the parameters into field domains, erases recognized proof fields, and checks the retained field layouts and metadata counts.  A nonrecursive inductive has no indices, at least one constructor, and a recognized retained layout for every constructor.  The predicates \longcode{structureFieldKindsWithParams?} and \code{variantLayoutWithParams?} implement those checks~\cite{types}.

A recursive family uses one pointer slot per value.  Its members must agree on family membership and parameter count, have no indices, and use the same concrete parameters at recursive occurrences.  The recursive-field recognizer descends through arrays, products, \code{PSum}, \code{Option}, and \code{Except}.  Other fields use its explicit fallback to ordinary type recognition.  The exact fallback controls compound fields involving other recursive declarations.  A recognized recursive representation still needs an accepted recursor shape for traversal.

\subsection{Internal and public positions}

Let $\I$ be the internal-layout domain, $\Ap$ the public array-element domain, and $\Pp$ the public parameter and result domain.  These are inductively defined by Table~\ref{tab:domains}.  All quantification over fields includes inactive variant constructors.  Empty retained field lists satisfy the field premise.

\begin{table}[ht]
\centering\small
\begin{tabular}{L{0.28\linewidth}L{0.22\linewidth}L{0.19\linewidth}L{0.19\linewidth}}
\toprule
Type & $\I$ & $\Ap$ & $\Pp$\\
\midrule
$\ty{Unit}$ & Included & Excluded & Excluded\\
$\ty{Bool}$, $\ty{U8}$, $\ty{U32}$, $\ty{U64}$, $\ty{Nat64}$, $\ty{Bytes}$ & Included & Included & Included\\
$\ty{Array}\;\tau$ & $\tau\in\I$ & $\tau\in\Ap$ & $\tau\in\Ap$\\
$\tau\times\upsilon$, $\ty{PSum}\;\tau\;\upsilon$ & Components in $\I$ & Excluded & Excluded\\
Structure or variant & Fields in $\I$ & Fields in $\Ap$ & Fields in $\Pp$\\
$\ty{Rec}(D,\bar\tau)$ & Included & Excluded & Excluded\\
\bottomrule
\end{tabular}
\caption{Position-sensitive runtime type domains.}
\label{tab:domains}
\end{table}

The implementation predicates are \code{valueLayout?}, \code{arrayElementLayout?}, \longcode{supportedPublicArrayElementType}, \code{supportedParamAbiType}, and \code{supportedResultAbiType}~\cite{types}.  Internal array-element layout and internal value layout admit the same type constructors.  Public arrays contain internal element layouts in memory, even though the public array value crosses the entry interface as one pointer.

\begin{proposition}[Public-domain inclusion]
For the domains in Table~\ref{tab:domains}, $\Pp=\Ap\subseteq\I$.  A type in $\Pp$ contains no product, \code{PSum}, recursive-pointer type, or runtime \code{Unit} at any retained field or array-element position.
\end{proposition}
\begin{proof}
Use induction on the finite type syntax.  The scalar and byte cases have identical public membership.  An array has membership in either public domain exactly when its element belongs to $\Ap$.  For structures and variants, the induction hypotheses equate the public membership of each field.  The other constructors belong to neither public domain.  This proves equality.  For inclusion in $\I$, the scalar and byte cases hold by definition, the array case follows from its element induction hypothesis, and the field cases follow componentwise.  The exclusion statement follows from the same induction because no public introduction rule admits an excluded constructor.
\end{proof}

This proposition concerns the stated domains.  Establishing that a mechanized version of these rules agrees with the executable predicates remains part of the implementation-equivalence work.

\subsection{Slot widths}

Write $w_I(\tau)$ and $w_P(\tau)$ for internal and public slot counts.  Each scalar has width one.  Internal \code{Unit} and recursive pointers have width one.  Bytes have internal width three and public width two.  Arrays have internal width two and public width one.  Products concatenate internal layouts, and \code{PSum} adds one tag slot.  For $X\in\{I,P\}$,
\begin{align*}
w_X(\ty{Struct}(S,\bar\tau,fs))&=\sum_{\upsilon\in fs}w_X(\upsilon),\\
w_X(\ty{Variant}(D,\bar\tau,cs))&=1+\sum_{fs\in cs}\sum_{\upsilon\in fs}w_X(\upsilon).
\end{align*}
These equations reserve payload slots for every variant constructor.  Public widths apply only on the public domain.  The total implementation function \code{abiSlots} returning a number for another type supplies no public-layout derivation~\cite{types}.

Define the public owner-slot count $h$ on $\Pp$ by $h=0$ for scalars, $h(\ty{Bytes})=h(\ty{Array}\;\tau)=1$, and by summing $h$ over structure fields or all variant payload fields.  Array contents are excluded from that sum because their element owner slots remain in memory in both entry conventions.

\begin{proposition}[Owner-slot difference]
For every $\tau\in\Pp$, $w_I(\tau)=w_P(\tau)+h(\tau)$.
\end{proposition}
\begin{proof}
Induct on the public-type derivation.  The scalar cases give $1=1+0$.  Bytes give $3=2+1$, and arrays give $2=1+1$.  A structure follows by summing the induction hypotheses over its fields.  A variant follows by summing over all constructor fields and adding the same tag slot on both sides.
\end{proof}

A structure containing one byte array and one array has internal width five, public width three, and owner-slot count two.  Its optional form has widths six and four.  These are consequences of the definitions and describe precisely which top-level slots the entry adapter omits.

\section{First-order terms and extraction}

\subsection{Runtime typing rules}

A runtime context is well formed when every distinct variable has a type in $\I$.  A signature context assigns each distinct function name a finite parameter list and result in $\I$.  An exported signature requires those types in $\Pp$.  The judgment $\Xi;\Delta\vdash e:\tau$ describes the abstract runtime term presentation.

Terms include variables, scalar literals, constructors, projection, let binding, branches, matches, direct calls, typed primitive operations, traps, and code bodies for recognized iteration and recursion.  A code body has explicit local binders that extraction substitutes or turns into first-order helper parameters.  Runtime arrow types and function-valued fields have no formation rule.

The variable and direct-call rules are
\[
\frac{x:\tau\in\Delta}{\Xi;\Delta\vdash x:\tau}
\qquad
\frac{\Xi(f)=(\tau_1,\ldots,\tau_n)\to\tau\quad
      \Xi;\Delta\vdash a_i:\tau_i\;(1\le i\le n)}
     {\Xi;\Delta\vdash f(\bar a):\tau}.
\]
Let and branch typing require
\[
\frac{\Xi;\Delta\vdash a:\tau\quad\Xi;\Delta,x:\tau\vdash b:\upsilon}
     {\Xi;\Delta\vdash\ty{let}\;x=a\;\ty{in}\;b:\upsilon}
\]
and
\[
\frac{\Xi;\Delta\vdash c:\ty{Bool}\quad
      \Xi;\Delta\vdash t:\tau\quad\Xi;\Delta\vdash u:\tau}
     {\Xi;\Delta\vdash\ty{if}\;c\;\ty{then}\;t\;\ty{else}\;u:\tau}.
\]
A structure constructor checks its ordered retained field types, and projection returns the designated field type.  A constructor $k$ of a recognized variant or recursive family $D\bar\tau$ checks its active retained field list $[\upsilon_{k,1},\ldots,\upsilon_{k,n_k}]$.  Its match rule is
\[
\frac{\Xi;\Delta\vdash a:D\bar\tau\quad\rho\in\I\quad
 \bigl(\Xi;\Delta,\bar x_k:\bar\upsilon_k\vdash e_k:\rho\bigr)_{k}}
 {\Xi;\Delta\vdash\ty{match}\;a\;\ty{with}\;\{k(\bar x_k)\mapsto e_k\}_k:\rho}.
\]
Each arm has fresh local pattern variables.  Products use pair construction and the two corresponding projection rules.  A nonrecursive Nat match extends only its successor arm with a \code{Nat64} predecessor.  A trap may have any represented result type and denotes terminal machine failure when evaluated.

These typing rules state necessary runtime shape conditions.  Accepted sparse matches, dependent conditionals, and generated recursors also require their exact implemented recognizers.  The abstract rules alone therefore do not decide whether LeanExe accepts a checked Lean expression.

\subsection{Primitive and callback signatures}

Let $\ty{Num}=\{\ty{U8},\ty{U32},\ty{U64},\ty{Nat64}\}$.  Structural equality applies to scalars and bytes, and to arrays, products, sums, structures, and nonrecursive variants whose components admit equality.  A custom \code{BEq} instance requires specialization to its selected body.  Recursive-inductive equality has no structural-equality case~\cite{demand,extract}.

\begin{longtable}{L{0.34\linewidth}L{0.59\linewidth}}
\caption{Runtime primitive signatures after static-argument removal.}\label{tab:primitives}\\
\toprule Family & Signature or code-body requirement\\\midrule\endfirsthead
\toprule Family & Signature or code-body requirement\\\midrule\endhead
Numeric arithmetic & $\nu\times\nu\to\nu$, $\nu\in\ty{Num}$, for add, subtract, multiply, divide, remainder, minimum, and maximum.\\
Numeric comparison & $\nu\times\nu\to\ty{Bool}$ for unsigned comparisons.\\
Fixed-width bits & $\ty{Uw}\times\ty{Uw}\to\ty{Uw}$ for bitwise operations and shifts; $\ty{Uw}\to\ty{Uw}$ for complement, $w\in\{8,32,64\}$.\\
Numeric conversion & Concrete conversions between distinct members of $\ty{Num}$; also $\ty{Bool}\to\ty{Nat64}$.\\
Binary64 intrinsics & Four operations $\ty{U64}\times\ty{U64}\to\ty{U64}$ and square root $\ty{U64}\to\ty{U64}$.\\
Array reads & Size returns \code{Nat64}; emptiness returns \code{Bool}; indexed read returns $\tau$ or $\ty{Option}\;\tau$; default read also takes a fallback $\tau$.\\
Array updates & Set and insert take an array, index, and element.  Erase takes an index.  Push, pop, append, extract, swap, and reverse return an array of the same element type.\\
Array map and modify & Body context extends $\Delta$ by $x:\tau$.  Map body has type $\upsilon$ and returns $\ty{Array}\;\upsilon$; modify body has type $\tau$.\\
Array predicates & A body $x:\tau\vdash p:\ty{Bool}$ supplies filter, find, optional find-index, any, and all.\\
Byte operations & Sequence operations use \code{Bytes}, \code{U8} elements, and \code{Nat64} indices.  Eight-byte endian decoding returns \code{U64}.\\
Pure folds & A left-fold body extends $\Delta$ by $a:\alpha,x:\tau$ and returns $\alpha$.  Right array folds bind element then accumulator.  Byte folds use $\tau=\ty{U8}$.\\
Monadic folds & The left-fold body returns $F\alpha$ for $F=\ty{Option}$ or $F=\ty{Except}\;\epsilon$.\\
Option and Except & Constructors, tag tests, maps, binds, defaults, elimination, and accepted fallback forms have their specialized payload types.  Bind bodies return the same selected monadic family.\\
Runtime intrinsics & Counter reads return \code{U64}.  Explicit release takes an accepted array or recursive root and returns \code{U64}, subject to the ownership judgment.\\
\bottomrule
\end{longtable}

For a code-bearing operation with ordinary argument types $\bar\tau$, body variables $\bar x:\bar\upsilon$, body result $\alpha$, and final result $\beta$, the typing premise is
\[
\frac{(\Xi;\Delta\vdash a_i:\tau_i)_i\quad
       \Xi;\Delta,\bar x:\bar\upsilon\vdash b:\alpha}
      {\Xi;\Delta\vdash p(\bar a;\bar x\mapsto b):\beta}.
\]
The body can refer to existing variables in $\Delta$.  Specialization must preserve their binding or pass them as explicit captured parameters.  Optional iteration bounds add \code{Nat64} arguments.  \code{extractPrimitivePairFrom}, \code{extractValueFrom}, and the pattern helpers fix the accepted source names and argument shapes~\cite{extract,patterns}.

\subsection{Static specialization and demand}

For a helper type $(x_1:A_1)\to\cdots\to(x_n:A_n)\to B$ applied to $\bar a$, \code{specializedInlineCall?} processes binders in source order.  It substitutes preceding arguments into each domain and runs \code{betaReduceExpr} with fuel 32.  A domain recognized in $\I$ contributes a runtime binder.  Otherwise, a successful \code{staticInlineArg} test contributes a static input.  Static inputs include sorts, recognized proofs, class evidence, and direct lambdas in function domains.  The substituted result must have a supported runtime type~\cite{types}.

\code{specializeInlineValue} substitutes static inputs and retains runtime binders.  Bounded normalization reduces class-method projections and selected transparent applications.  \longcode{blocksTransparentSpecialization} protects the primitive, matcher, recursor, and monad families that have dedicated lowering.  A generic source helper can therefore have accepted concrete inline uses while its unspecialized type has no retained runtime signature.

Demand analysis controls strict materialization of otherwise deferred values.  Write $\operatorname{MayTrap}(a)$ for the executable possible-trap analysis and $\operatorname{Must}(f,i)$ for its necessary-demand summary.  The tests include
\begin{align*}
\operatorname{StrictCallSafe}(f,\bar a)
  &=\bigwedge_i(\neg\operatorname{MayTrap}(a_i)\lor\operatorname{Must}(f,i)),\\
\operatorname{MaterializeSafe}(\bar\tau,\bar a)
  &=\bigwedge_i(w_I(\tau_i)=1\lor\neg\operatorname{MayTrap}(a_i)).
\end{align*}
These equations describe the implemented analyses~\cite{demand}.  Their interpretation as sound semantic summaries is an open theorem.  Lazy bindings can preserve an unused field or branch whose evaluation would trap.  Strict helper arguments and fixed-width array elements stage multi-slot results in locals, so projections reuse one evaluated call result~\cite{storage,extract}.

\subsection{Recognized control and recursion}

The accepted control forms include \code{Id}, \code{Option}, and concrete \code{Except} sequencing, direct matches, and \code{ForIn.forIn} over byte arrays, fixed-width arrays, legacy ranges, and \code{Lean.Loop}.  The loop accumulator must have a supported internal layout.  Generated \code{ForInStep.done}, failure tags, and accepted branch exits determine termination of the lowered loop~\cite{patterns,extract}.

Recursion recognition has six relevant families: first-\code{Nat} fuel recursion, direct structural recursion, expression-position structural recursion with synthetic helpers, closed list-shaped folds and predicates, array descent through the recognized \code{WellFounded.fix} shape, and mutual descent through a nested \code{PSum} form of \code{WellFounded.Nat.fix}.  Direct structural descent uses recognized generated below projections.  Array descent uses attached elements and the generated recursive handle.  The mutual form immediately dispatches to a supported family member~\cite{structural,extract}.  Lean termination checking precedes these shape restrictions.  Complete independent inference rules for these generated forms remain part of the formalization work.

\section{Value representation and operational rules}

\subsection{Memory and ownership}

An abstract byte array is a finite byte sequence, an array is a finite sequence of values of its element type, and a recursive value is a finite constructor tree.  A physical representation may share tree subvalues.  Let $\mu$ be byte-addressed linear memory and $\operatorname{load}_{64}(\mu,p)$ a valid little-endian word load.  Internal bytes have slots $[o,p,n]$, where $o$ is an owner root, $p$ the visible byte pointer, and $n$ the length.  Public bytes expose $[p,n]$.  Internal arrays have $[o,p]$, and public arrays expose $[p]$.

For bytes $b$, the payload relation requires $n=|b|$ and $\mu[p:p+n]=b$.  Owner zero permits borrowed storage.  A nonzero owner requires a live allocation root whose lifetime covers the visible range.  For an array $a:\ty{Array}\;\tau$, the header and elements satisfy
\[
\operatorname{load}_{64}(\mu,p)=|a|,
\qquad
a_i\text{ is represented at }p+8+8i\,w_I(\tau).
\]
Each element uses its internal layout, including embedded owner slots.  Structures concatenate field vectors.  A variant stores its constructor index and all reserved payload areas.  A recursive value uses a pointer to a tagged slot object~\cite{storage,values,binary}.

The representation relation includes valid ranges, live owners, and alias accounting appropriate to the performed operations.  These are program-theorem premises.  A general representation-preservation theorem must justify them for every accepted operation.  Inactive scalar payloads have no source value.  Inactive owner positions still participate in a type-derived child-mask walk if retained in a container, so their contents require explicit heap premises.  A public byte-array result also omits the owner identity needed to distinguish an allocation root from a slice.

All address observations require bounds that prevent relevant machine arithmetic from wrapping and place each access in linear memory.  The emitted code computes addresses with \code{i64} arithmetic and converts them with \code{i32.wrap\_i64}.  The complete target relation retains this machine behavior on every input.  A source-facing theorem states the narrower represented-input and resource premises under which the natural-number address formula applies.

\subsection{Numeric semantics}

Let $U=2^{64}$.  Fixed-width unsigned addition, subtraction, and multiplication reduce modulo $2^w$.  Narrow operations and conversions normalize the result to width $w\in\{8,32\}$.  Unsigned division returns zero for a zero divisor, and remainder returns the dividend in that case.  Shift counts reduce modulo the operand width.  Comparisons use unsigned order~\cite{extract,binary}.

Runtime natural numbers lie in $[0,U)$.  Their subtraction saturates at zero.  Addition and multiplication return the mathematical result when it lies below $U$ and trap otherwise.  Successor uses checked addition, and predecessor uses saturating subtraction.  The backend implements the addition test by detecting wrapped order, and the multiplication test by comparing against $\lfloor(U-1)/b\rfloor$ after handling $b=0$~\cite{binary}.

The five binary64 operations retain raw \code{U64} arguments and results.  Write $a_f=\operatorname{reinterpret}_{f64}(a)$ and $b_f=\operatorname{reinterpret}_{f64}(b)$.  For $p\in\{+,-,\times,/\}$, let $\operatorname{Wasm}_{p,64}(a_f,b_f)$ be the set of results permitted by the WebAssembly binary64 operation.  Define
\[
\operatorname{Fp}_p(a,b,r)
\iff
r\in\{\operatorname{reinterpret}_{i64}(z)\mid
 z\in\operatorname{Wasm}_{p,64}(a_f,b_f)\}.
\]
Square root has the corresponding unary relation.  WebAssembly permits a set of NaN results in applicable cases~\cite{wasm}.  A statement over all conforming engines must account for that relation.  Exact-word theorems in the pinned Talos model identify that model's arithmetic result.  The source \code{LeanExe.Float64} definitions provide native Lean comparison functions and require a separate bridge for an across-engine claim.

\subsection{Collections, branches, and loops}

Branches execute the selected arm, Boolean conjunction and disjunction short-circuit, and \code{Option}/\code{Except} sequencing skips continuations on failure.  The emitted instruction order fixes the schedule after demand-based extraction.  Direct calls evaluate materialized arguments in order and bind their result slots for reuse~\cite{extract,binary}.

Array operations observe finite-sequence behavior under valid representation and allocation premises.  Reads return an indexed element or an optional/default result.  Bang wrappers trap on invalid indices.  Set, insert, erase, and swap have their recognized conditional or trapping boundary.  Push, pop, append, extract, reverse, map, and filter preserve the stated sequence ordering.  Updated arrays allocate and copy as specified by their emitter.  Unchanged conditional results preserve their owner.  Copying existing child roots retains them, while recognized fresh children can transfer ownership~\cite{values,binary}.

Find operations stop at the first true predicate result.  Any and all stop at their first decisive result.  Equality checks length before its ordered scan.  Forward array and byte folds visit $[\mathit{start},\min(\mathit{stop},n))$.  Reverse array folds begin at $\min(\mathit{start},n)$ and visit the preceding element while that index exceeds the exclusive lower bound.  Range loops increment through checked natural-number addition.  A zero step reaching that loop repeats the index until a body exit.  Monadic folds and loops stop on their first failure tag~\cite{binary}.

Accumulator replacement stages all new slots before releasing selected previous roots and copying the staged slots.  The first accumulator value is excluded from this release rule because source aliases may retain it.  Later release selection depends on fresh-owner and liveness analyses.  The emitted multi-slot assignment retains a shared fold computation when any result slot remains live~\cite{values,binary}.

Byte slicing preserves the owner and adjusts the visible range.  Endian conversion to \code{UInt64} requires exactly eight bytes, returning $\sum_i b_i2^{8i}$ for little endian or $\sum_i b_i2^{8(7-i)}$ for big endian.  The byte-copy-slice operation forms the preserved destination prefix, the available requested source bytes, and the remaining destination suffix, with the implemented clamping rules~\cite{values,binary}.  Compile-time string recognition accepts its specified ASCII forms.  Runtime ASCII and restricted JSON helpers are source libraries over the same bytes and recursive-data rules~\cite{specification}.

\section{Compilation and target execution}

\subsection{The exact compilation graph}

For a partial implementation function $g$ at revision $v$, write $G_v(g,x,y)$ when evaluating $g$ on $x$ terminates with $y$.  This relation fixes successful and rejected results without asserting totality of every compiler function.  Let $Q=(E,m,f,\mathit{mode},\mathit{bounds})$ identify the checked environment, module, entry, compilation mode, and mode bounds.

For library extraction,
\[
\operatorname{Extract}_v(E,m,f,\ty{library},M)
\iff G_v(\code{compileEnvironment},(E,m,f),\ok(M)).
\]
The fixed command modes use their corresponding environment-extraction functions.  \longcode{compileEnvironmentWithEntryModeDetailed} checks the entry, collects the reachable namespace-rooted declarations, normalizes with its fixed fuel, and collects synthetic helpers.  It extracts once with empty fresh-result summaries, derives those summaries, validates explicit releases, and extracts again using the summaries~\cite{extract}.  Report classification and complete extraction therefore have different result objects.

Let $\operatorname{Emit}_v(\mathit{mode},\mathit{bounds},f,M)$ name the selected production serializer.  Treat the library serializer's direct byte result as an $\ok$ result for uniform notation.  Then
\begin{align*}
\operatorname{Compile}_v(Q,b)\iff\exists M.\;&
\operatorname{Extract}_v(E,m,f,\mathit{mode},M)\\
&\land\operatorname{Emit}_v(\mathit{mode},\mathit{bounds},f,M)=\ok(b).
\end{align*}
The serializer lowers IR functions to structured instruction lists and assembles the binary.  Library \code{moduleBytes} calls \code{legacyModuleBytes}.  The image-emitter experiment has a separate validation path~\cite{binary}.  This distinction is relevant to the export-name concern below.

Define a separate executable-module judgment
\[
\operatorname{Executable}_v(Q,b,W)
\iff\operatorname{Compile}_v(Q,b)\land\operatorname{Decode}(b)=W
       \land\operatorname{Valid}(W).
\]
Here $\operatorname{Decode}$ and $\operatorname{Valid}$ belong to the chosen WebAssembly syntax and validation account.  Separating byte production from validity preserves the possibility of an accepted source that exposes a compiler defect.  WebAssembly stack typing checks target value types and control shapes~\cite{wasm}.  An \code{i64} stack slot alone supplies no source Boolean, narrow-integer, or pointer invariant.

\subsection{Allocator and fixed entry adapters}

Library mode exports memory, the selected entry, allocation and reference-count operations, and four counters.  Memory starts at 16 pages, and the heap starts at byte offset 4096.  Each owned root $p$ has a 48-byte header containing magic value, reference count, capacity, kind, width or slot count, and child mask.  The allocator rounds payloads to eight-byte alignment with a minimum of eight bytes, takes the first sufficiently large free block, or advances the heap and grows memory~\cite{binary}.

Retaining root zero changes no state.  Retaining a nonzero root checks magic and positive count, increments the retain counter, and increments the stored count.  Release of zero is also a no-op.  Release of a valid nonzero root increments its counter, decrements a count above one, or recursively releases marked children before adding the root to the free list.  Invalid headers and zero-count reuse trap.  All counters and reference-count increments use word arithmetic.  A natural-number accounting theorem therefore needs suitable overflow premises.

Explicit source release requires \code{validateModuleReleases}.  An accepted release consumes a direct fresh allocation, a helper root identified as fresh, or a statically owner-zero array at its supported final-use position~\cite{release}.  Automatic cleanup additionally uses result-root, borrowing, and liveness analyses.  The source definitions of the runtime counters and release all return zero, while generated execution reads and updates allocator state.  A source refinement statement that observes those operations needs their extended semantics.

The exported reset function restores the heap top, free list, and counters to their initial values.  It leaves the page count and old payload bytes in memory.  Earlier pointers lose their lifetime guarantee because allocation can reuse those bytes.  A public byte result's pointer and length should therefore be interpreted with its program-specific lifetime premise.

The five WASI command shapes accept a zero-argument byte result, a bounded stdin byte transform, a stdin transform returning \code{Except}, an argv transform returning \code{Except}, or the combined stdin/argv transform.  The pure entry retains the corresponding source type.  The generated adapter supplies the fixed imports and output trace~\cite{binary,specification}.

For maximum stdin bytes $b$, maximum user arguments $a$, and maximum argv-buffer bytes $g$, the source constants give
\[
b\le16\cdot65536-4096-1,
\qquad R(a,g)=8+24a+4(a+1)+g.
\]
Argv mode requires $R(a,g)\le16\cdot65536-4096$, and combined mode requires $b+8+R(a,g)\le16\cdot65536-4096$ as well as the stdin bound.  Runtime readers enforce their configured counts, skip \code{argv[0]}, and pass borrowed-owner values.  A successful result writes stdout.  An error result writes stderr and exits with status one.  Invalid tags, failed imports, exceeded bounds, and short writes trap.  These statements assume WASI imports implement their specified interface.

\subsection{Execution and observations}

Let $\longrightarrow_W$ be the selected WebAssembly small-step relation with memory, globals, locals, stack, frames, and a host-import relation.  Let $\operatorname{Init}(W,H,\bar a,c_0)$ describe an instantiated module and exported call in host state $H$.  An observation includes return slots and relevant memory, a trap, or a host exit and byte-output trace.  Define
\begin{align*}
\operatorname{Run}_v(Q,H,\bar a,o)\iff\exists b,W,c_0,c_n.\;&
\operatorname{Executable}_v(Q,b,W)\\
&\land\operatorname{Init}(W,H,\bar a,c_0)\\
&\land c_0\longrightarrow_W^*c_n
\land\operatorname{Observe}(c_n)=o.
\end{align*}
Infinite reduction defines divergence.  An external process limit is a separate operational event.  This relation covers all emitted instructions through the selected target semantics, including wrapping addresses, floating-point special values, and malformed host inputs.  Only inputs satisfying a stated representation relation acquire source-value interpretations.

The repository's program proofs instantiate execution with the pinned Talos model.  A claim about every conforming engine additionally requires an appropriate relation to the WebAssembly standard, or a theorem stated against that standard.  Conformance and runtime tests provide finite evidence for that relation.  Their results do not expand a theorem's quantified scope.

\section{Source audit and existing proof boundaries}
\label{sec:audit}

\subsection{Export names and child-mask width}

The source reserves only \code{memory}, \code{alloc}, and \code{reset} in \longcode{Extract.Core.reservedExportNames}.  Production export assembly also introduces \code{retain}, \code{release}, \code{free}, \code{allocCount}, \code{retainCount}, \code{releaseCount}, and \code{freeCount}~\cite{extract,binary}.  The prose specification reserves all ten names.  The production serializer bypasses the image-emitter duplicate-name check.  These declarations expose a gap between the documented rejection set and its implemented entry-name check.  A general emitted-validity theorem must resolve that gap.  Report preparation changed no acceptance rule and ran no example named after an omitted export.

The ownership-mask boundary follows from three declarations.  \code{heapChildMaskFromType} constructs a natural-number mask with a set bit at each owner position.  \code{Binary.i64Const} reduces a mask constant modulo $2^{64}$.  \code{coreReleaseInstrs} scans object slots and tests the stored mask through \code{i64.shr\_u}, whose shift count is modulo 64~\cite{values,binary,wasm}.  The inspected layout predicates provide no general 64-slot bound.  Consequently a general ownership proof needs a bound connecting the accepted layout to mask representation, or a verified wider representation.  This is a source-level proof obligation and potential implementation defect.  A compiled reproducer remains open.

\subsection{Reference IR evaluation}

The IR reference evaluator has a restricted comparison interpretation.  \code{IR.Expr.eval} returns zero for traps, runtime counter reads, releases, heap allocation, heap loads, and collection cases that need memory.  Some update cases return an existing pointer expression.  Its \code{natAdd} and \code{natMul} use word arithmetic without the backend's overflow checks.  Its loops have a fixed evaluation bound~\cite{ir}.

\code{Extract.Eval.evalEntry} filters calls to scalar arguments, scalar-slot results, and a module passing \code{scalarModule}~\cite{ireval}.  That filter excludes heap operations but admits traps and checked-Nat arithmetic.  Thus equality with \code{eval-ir} supports a comparison only when the execution avoids those differences and stays within the evaluator bound.  The compiled execution relation above instead uses production lowering and target execution.

\subsection{Named theorem scope}

The existing scalar certificates connect successful descriptor recognition to exact structured emission.  For example, \longcode{ScalarDescriptor.Expr.ofIR\_emitWithRelease} assumes successful expression reification and proves equality with the backend expression emitter.  The condition, statement, loop, and encoded-index certificates make corresponding statements~\cite{certificates}.  These theorems establish emitter equalities at the named boundary.

The exact-artifact path embeds bytes, decodes them, validates the restricted module, identifies its Talos translation, and proves the package's behavioral proposition.  \longcode{Wasm.Binary.Proof.decode\_sound} proves that successful decoding satisfies \code{Grammar.Encodes}.  \longcode{Wasm.Binary.Proof.validate\_sound} proves that successful validation satisfies \code{CoreValid}~\cite{artifact}.  A package theorem identifies the translated decoded module with the model used by its behavioral proof.  The artifact format also records the external-file comparison and the proof's stated host assumptions.

The artifact profile excludes WASI imports, so its current decoder and validator theorems give no automatic certificate for the command adapters in this report.  Runtime helper theorems, including recursive teardown of represented ownership trees, have their own memory and separation premises~\cite{artifact}.  The compiler's source and the extraction passes remain outside an artifact-only theorem's assumptions.  Their agreement with the source program requires additional refinement evidence.

\section{Independent formalization and refinement obligations}

An independent formalization must replace the remaining recognizer references with declared syntax and judgments and prove their agreement with the implementation.  The abstract runtime typing rules, domain definitions, and layout propositions above supply part of that specification.  The shape-sensitive recursion rules, demand semantics, and ownership relation need complete independent definitions before a general preservation theorem can be stated precisely.

Extraction typing would prove that successful extraction produces a well-formed slot program with valid local indices, direct-call indices, parameter and result arities, and fixed-width offsets.  A stronger invariant would connect those slots to represented source values and live owner roots.  The present \code{IR.Expr} and \code{IR.Stmt} datatypes have no source-type indices that establish those properties by construction~\cite{ir}.

The remaining semantic obligations concern erasure, static substitution, demand analysis, source-operation lowering, ownership transfer, emitted validity, and serialization.  A source-to-target theorem must state the selected source evaluation relation, trap and divergence behavior, bounded arithmetic, representation, resources, and runtime-intrinsic observations.  For a nondeterministic target relation, it must constrain every allowed target outcome covered by the claim.  A source termination proof also needs the lowering relation before it establishes compiled termination.

A reference-evaluator comparison concerns its stated input and evaluator domain.  An exact-artifact theorem concerns its embedded binary and behavioral premises.  The implemented acceptance graph, the target execution relation, and the named artifact proofs each identify a concrete subject.  A complete compiler-correctness result requires proved connections between those subjects under explicit premises.

\begin{thebibliography}{99}
\small
\bibitem{checkpoint} LeanExe contributors. \emph{LeanExe source checkpoint}. Commit 2f3ec33f6e98ad98ac94df277c91e5319763cf9d, 2026. \url{https://github.com/jsmorph/leanexe/tree/2f3ec33f6e98ad98ac94df277c91e5319763cf9d}.
\bibitem{typetheory} LeanExe contributors. \emph{Type Theory of the LeanExe Fragment}. \url{https://github.com/jsmorph/leanexe/blob/2f3ec33f6e98ad98ac94df277c91e5319763cf9d/docs/leanexe-type-theory.md}.
\bibitem{specification} LeanExe contributors. \emph{Formal Specification of LeanExe Compilation and Execution}. \url{https://github.com/jsmorph/leanexe/blob/2f3ec33f6e98ad98ac94df277c91e5319763cf9d/docs/leanexe-formal-specification.md}.
\bibitem{lean} Lean developers. \emph{The Lean Language Reference: The Type System}. \url{https://lean-lang.org/doc/reference/latest/The-Type-System/}. Implementation pin in the cited LeanExe checkpoint: Lean 4.34.0-rc2, commit 6a10ac8c22beadecabdbb0919c2b50214762f91d.
\bibitem{types} LeanExe contributors. \emph{Type and layout recognition}. \url{https://github.com/jsmorph/leanexe/blob/2f3ec33f6e98ad98ac94df277c91e5319763cf9d/LeanExe/Extract/Types.lean}.
\bibitem{extract} LeanExe contributors. \emph{Expression and module extraction}. \url{https://github.com/jsmorph/leanexe/blob/2f3ec33f6e98ad98ac94df277c91e5319763cf9d/LeanExe/Extract/Core.lean}.
\bibitem{patterns} LeanExe contributors. \emph{Pattern and monad recognition}. \url{https://github.com/jsmorph/leanexe/blob/2f3ec33f6e98ad98ac94df277c91e5319763cf9d/LeanExe/Extract/Patterns.lean}.
\bibitem{structural} LeanExe contributors. \emph{Structural recursion recognition}. \url{https://github.com/jsmorph/leanexe/blob/2f3ec33f6e98ad98ac94df277c91e5319763cf9d/LeanExe/Extract/StructuralRec.lean}.
\bibitem{demand} LeanExe contributors. \emph{Demand analysis and structural equality}. \url{https://github.com/jsmorph/leanexe/blob/2f3ec33f6e98ad98ac94df277c91e5319763cf9d/LeanExe/Extract/Demand.lean}.
\bibitem{storage} LeanExe contributors. \emph{Storage and parameter lowering}. \url{https://github.com/jsmorph/leanexe/blob/2f3ec33f6e98ad98ac94df277c91e5319763cf9d/LeanExe/Extract/Storage.lean}.
\bibitem{values} LeanExe contributors. \emph{Value layouts, owner masks, and liveness}. \url{https://github.com/jsmorph/leanexe/blob/2f3ec33f6e98ad98ac94df277c91e5319763cf9d/LeanExe/Extract/Values.lean}.
\bibitem{release} LeanExe contributors. \emph{Explicit release checking}. \url{https://github.com/jsmorph/leanexe/blob/2f3ec33f6e98ad98ac94df277c91e5319763cf9d/LeanExe/Extract/ReleaseCheck.lean}.
\bibitem{binary} LeanExe contributors. \emph{Production WebAssembly lowering and binary emission}. \url{https://github.com/jsmorph/leanexe/blob/2f3ec33f6e98ad98ac94df277c91e5319763cf9d/LeanExe/Wasm/Binary.lean}.
\bibitem{ir} LeanExe contributors. \emph{IR syntax and reference evaluation}. \url{https://github.com/jsmorph/leanexe/blob/2f3ec33f6e98ad98ac94df277c91e5319763cf9d/LeanExe/IR/Core.lean}.
\bibitem{ireval} LeanExe contributors. \emph{Reference-evaluator entry restriction}. \url{https://github.com/jsmorph/leanexe/blob/2f3ec33f6e98ad98ac94df277c91e5319763cf9d/LeanExe/Extract/Eval.lean}.
\bibitem{wasm} WebAssembly Community Group. \emph{WebAssembly Core Specification}. Numeric semantics and validation. \url{https://webassembly.github.io/spec/core/exec/numerics.html}; \url{https://webassembly.github.io/spec/core/valid/instructions.html}. Accessed 11 September 2026.
\bibitem{certificates} LeanExe contributors. \emph{Scalar emitter certificates}. \url{https://github.com/jsmorph/leanexe/blob/2f3ec33f6e98ad98ac94df277c91e5319763cf9d/LeanExe/Wasm/ScalarCertificate.lean}.
\bibitem{artifact} LeanExe contributors. \emph{Artifact Verification Format and Talos proof inventory}. \url{https://github.com/jsmorph/leanexe/blob/2f3ec33f6e98ad98ac94df277c91e5319763cf9d/docs/artifact-format.md}; \url{https://github.com/jsmorph/leanexe/blob/2f3ec33f6e98ad98ac94df277c91e5319763cf9d/proofs/talos/README.md}. Theorem sources: \longcode{Project/Artifact/Binary/Proof/Decode.lean} and \code{Validate.lean}, namespace \code{Wasm.Binary.Proof}.
\end{thebibliography}
\end{document}
